\documentclass[aspectratio=169,11pt,english]{beamer}
\input{../../../_preambulo_beamer.tex}
\input{../../../_comandos_beamer.tex}
\title[L05 · Foundations]{Set Theory and Sample Spaces}
\begin{document}
\begin{frame}{L05 · An experiment and its outcomes}
A \textbf{random experiment} has an outcome we do not know before it occurs.
Consider one roll of a six-sided die.
\[
S=\{1,2,3,4,5,6\},\qquad |S|=6.
\]
The \textbf{sample space} $S$ contains every possible outcome. One roll produces one element of $S$.

\medskip
The \textbf{cardinality} $|S|$ is the number of distinct elements of $S$.
\end{frame}
\begin{frame}{An event is a subset}
An \textbf{event} collects outcomes satisfying a condition. For the die, define
\[
A=\{2,4,6\}\quad\text{(an even number)},\qquad
B=\{4,5,6\}\quad\text{(at least four)}.
\]
Both are subsets of $S$, and $|A|=|B|=3$.

\medskip
If the result is 4, both $A$ and $B$ occur: one outcome can belong to more than one event.
\end{frame}
\begin{frame}{Intersection: satisfy both conditions}
The intersection keeps only shared elements:
\[
A\cap B=\{x\in S:x\in A\text{ and }x\in B\}.
\]
In this example,
\[
A\cap B=\{4,6\},\qquad |A\cap B|=2.
\]
Four and six are even and are also at least four.
\end{frame}
\begin{frame}{Union: satisfy at least one condition}
The union contains elements of $A$, of $B$, or of both:
\[
A\cup B=\{x\in S:x\in A\text{ or }x\in B\}.
\]
The \textbf{or} is inclusive. Therefore,
\[
A\cup B=\{2,4,5,6\},\qquad |A\cup B|=4.
\]
Each outcome is listed once in a set.
\end{frame}
\begin{frame}{Count a union without double-counting}
Adding $|A|+|B|$ counts each element of $A\cap B$ twice. Subtract one copy of each:
\[
|A\cup B|=|A|+|B|-|A\cap B|.
\]
For the events above:
\[
4=3+3-2.
\]
This identity holds for any two finite sets.
\end{frame}
\begin{frame}{Complement: relative to a universe}
The complement of $A$ contains outcomes in $S$ that are not in $A$:
\[
A^c=S\setminus A=\{1,3,5\},\qquad |A^c|=3.
\]
$A$ and $A^c$ partition the sample space:
\[
A\cap A^c=\varnothing,\quad A\cup A^c=S,\quad
|A^c|=|S|-|A|=6-3=3.
\]
A complement is ambiguous until its universe is specified.
\end{frame}
\begin{frame}{Disjointness and exhaustiveness}
Two events $X,Y\subseteq S$ are \textbf{disjoint} if $X\cap Y=\varnothing$.
They are \textbf{exhaustive} if $X\cup Y=S$.

\medskip
Let $C=\{1,3,5\}$ (odd) and $D=\{2,4,6\}$ (even). Then
\[
C\cap D=\varnothing,\qquad C\cup D=S.
\]
$C$ and $D$ have both properties. Disjointness concerns \textbf{overlap}; exhaustiveness concerns \textbf{coverage}.
\end{frame}
\begin{frame}{Checks for your results}
For events $A,B\subseteq S$:
\[
A\cap B\subseteq A,\quad A\cap B\subseteq B,\qquad
A\subseteq A\cup B,\quad B\subseteq A\cup B.
\]
Hence $|A\cap B|$ cannot exceed $|A|$ or $|B|$, and $|A\cup B|$ cannot be smaller than either one.

\medskip
\scriptsize Reference: OpenStax, \emph{Introductory Business Statistics 2e}, Sections 3.1 and 3.2. The die example is self-contained.
\end{frame}
\end{document}
